%====================================================%
% SEZIONE 1: Definizione
%====================================================%
\begin{question}[ID=1.1, subtitle={Calcolo di logaritmi con le proprietà}, difficulty=facile, tags={logaritmi, proprietà}]
Calcola il valore dei seguenti logaritmi usando le opportune proprietà:
\begin{center}\begin{minipage}{0.92\linewidth}
\begin{multicols}{2}
\begin{enumerate}[leftmargin=3.2cm]
\item $\log_2 8$
\item $\log 50 + \log 2$
\item $\log_3 \dfrac{x^3}{y} \color{gray}$
\item $\color{gray}\ln 4e^5 - \ln e^2 - 2 \ln 2$
\end{enumerate}
\end{multicols}
\end{minipage}\end{center}
\begin{solution}
\begin{enumerate}
    \item $3$
    \item $2$
    \item $3\log_3x-\log_3y$
    \item $3$
\end{enumerate}
\end{solution}
\begin{explanation}
Proprietà dei logaritmi (argomenti positivi): \ $\log_a(xy)=\log_ax+\log_ay$; \ $\log_a\dfrac xy=\log_ax-\log_ay$; \ $\log_a x^n=n\log_ax$; \ $\log_a a^n=n$.
\begin{enumerate}
    \item $\log_2 8=\log_2 2^3=3$.
    \item $\log50+\log2=\log(50\cdot2)=\log100=2$.
    \item $\log_3\dfrac{x^3}{y}=\log_3x^3-\log_3y=3\log_3x-\log_3y$ \ (per $x,y>0$).
    \item \[\begin{aligned}
    \ln4e^5-\ln e^2-2\ln2 &= \ln4+\ln e^5-2-\ln4 && \textcolor{gray}{\text{$2\ln2=\ln4$}}\\
    &= 5-2=3
    \end{aligned}\]
\end{enumerate}
\end{explanation}
\end{question}

%---------------------------------------------------------------------------------%

\begin{question}[ID=1.2, subtitle={Calcolo di logaritmi con le proprietà (variante)}, difficulty=facile, tags={logaritmi, proprietà}]
Calcola il valore dei seguenti logaritmi usando le opportune proprietà:
\begin{center}\begin{minipage}{0.92\linewidth}
\begin{multicols}{2}
\begin{enumerate}[leftmargin=3.2cm]
\item $\log_3 27$
\item $\log 20 + \log 5$
\item $\log_2 \left(\dfrac{x^4}{y^2}\right) \color{gray}$
\item $\color{gray}\ln 9e^3 - \ln e - 2\ln 3$
\end{enumerate}
\end{multicols}
\end{minipage}\end{center}
\begin{solution}
\begin{enumerate}
    \item $3$
    \item $2$
    \item $4\log_2x-2\log_2y$
    \item $2$
\end{enumerate}
\end{solution}
\begin{explanation}
Proprietà dei logaritmi (argomenti positivi): \ $\log_a(xy)=\log_ax+\log_ay$; \ $\log_a\dfrac xy=\log_ax-\log_ay$; \ $\log_a x^n=n\log_ax$; \ $\log_a a^n=n$.
\begin{enumerate}
    \item $\log_3 27=\log_3 3^3=3$.
    \item $\log20+\log5=\log100=2$.
    \item $\log_2x^4-\log_2y^2=4\log_2x-2\log_2y$.
    \item \[\begin{aligned}
    \ln9e^3-\ln e-2\ln3 &= \ln9+3-1-\ln9 && \textcolor{gray}{\text{$2\ln3=\ln9$}}\\
    &= 2
    \end{aligned}\]
\end{enumerate}
\end{explanation}
\end{question}


%====================================================%
% SEZIONE 2: Equazioni
%====================================================%
\begin{question}[ID=2.1, subtitle={Equazioni logaritmiche elementari}, difficulty=facile, tags={equazioni logaritmiche, definizione}]
Risolvi le seguenti equazioni logaritmiche:
\begin{center}\begin{minipage}{0.92\linewidth}
\begin{multicols}{2}
\begin{enumerate}[leftmargin=3.2cm]
\item $\log_2 x = 3$
\item $\log_5(2x-1) = 2$
\item $\log(2x+1) = \log(x-1)$
\item $\log_3 x = 0 \color{gray}$
\item $\color{gray}\log_3(x-4) + \log_3(x+4) = 2$
\item $4^x = 3$
\end{enumerate}
\end{multicols}
\end{minipage}\end{center}
\begin{solution}
\begin{enumerate}
    \item $x=8$
    \item $x=13$
    \item impossibile
    \item $x=1$
    \item $x=5$
    \item $x=\log_43$
\end{enumerate}
\end{solution}
\begin{explanation}
Prima si scrivono le \textbf{condizioni di esistenza} (argomenti positivi); poi si usa la definizione $\log_a A=b \iff A=a^b$, oppure il confronto $\log_aA=\log_aB \iff A=B$; infine si controllano le C.E.
\begin{enumerate}
    \item $x=2^3=8$.
    \item C.E. $x>\tfrac12$. \ $2x-1=5^2=25 \Rightarrow x=13$, accettabile.
    \item C.E. $2x+1>0$ e $x-1>0$, cioè $x>1$. \ Confronto: $2x+1=x-1 \Rightarrow x=-2$, che non rispetta le C.E.: impossibile.
    \item $x=3^0=1$.
    \item C.E. $x>4$.
    \[\begin{aligned}
    \log_3\big[(x-4)(x+4)\big] &= 2 && \textcolor{gray}{\text{somma di logaritmi}}\\
    x^2-16 &= 9 \;\Rightarrow\; x=\pm5
    \end{aligned}\]
    Solo $x=5$ rispetta le C.E.
    \item Applico il logaritmo in base $4$: \ $x=\log_43$.
\end{enumerate}
\end{explanation}
\end{question}

%---------------------------------------------------------------------------------%

\begin{question}[ID=2.2, subtitle={Equazioni logaritmiche elementari (variante)}, difficulty=facile, tags={equazioni logaritmiche, definizione}]
Risolvi le seguenti equazioni logaritmiche:
\begin{center}\begin{minipage}{0.92\linewidth}
\begin{multicols}{2}
\begin{enumerate}[leftmargin=3.2cm]
\item $\log_3 x = 2$
\item $\log_4(3x-2) = 1$
\item $\log(3x-1) = \log(x+5)$
\item $\log_2 x = -1 \color{gray}$
\item $\color{gray}\log_2(x-1) + \log_2(x+1) = 3$
\item $5^x = 7$
\end{enumerate}
\end{multicols}
\end{minipage}\end{center}
\begin{solution}
\begin{enumerate}
    \item $x=9$
    \item $x=2$
    \item $x=3$
    \item $x=\tfrac12$
    \item $x=3$
    \item $x=\log_57$
\end{enumerate}
\end{solution}
\begin{explanation}
Prima si scrivono le \textbf{condizioni di esistenza} (argomenti positivi); poi si usa la definizione $\log_a A=b \iff A=a^b$, oppure il confronto $\log_aA=\log_aB \iff A=B$; infine si controllano le C.E.
\begin{enumerate}
    \item $x=3^2=9$.
    \item C.E. $x>\tfrac23$. \ $3x-2=4 \Rightarrow x=2$.
    \item C.E. $x>\tfrac13$. \ $3x-1=x+5 \Rightarrow x=3$, accettabile.
    \item $x=2^{-1}=\tfrac12$.
    \item C.E. $x>1$. \ $\log_2(x^2-1)=3 \Rightarrow x^2-1=8 \Rightarrow x=\pm3$; solo $x=3$ accettabile.
    \item $x=\log_57$.
\end{enumerate}
\end{explanation}
\end{question}

%---------------------------------------------------------------------------------%

\begin{question}[ID=2.3, subtitle={Equazioni logaritmiche per confronto}, difficulty=media, tags={equazioni logaritmiche, confronto}]
Risolvi per confronto:
\begin{center}\begin{minipage}{0.92\linewidth}
\begin{multicols}{2}
\begin{enumerate}[leftmargin=3.2cm]
\item $\log_2(x) = \log_2(16)$
\item $\log_3(x+1) = \log_3(4)$
\item $\log(x-2) = \log(3)$
\item $\log_4(x-1) = \log_4(3)$
\end{enumerate}
\end{multicols}
\end{minipage}\end{center}
\begin{solution}
\begin{enumerate}
    \item $x=16$
    \item $x=3$
    \item $x=5$
    \item $x=4$
\end{enumerate}
\end{solution}
\begin{explanation}
Prima si scrivono le \textbf{condizioni di esistenza} (argomenti positivi); poi si usa la definizione $\log_a A=b \iff A=a^b$, oppure il confronto $\log_aA=\log_aB \iff A=B$; infine si controllano le C.E.
Stessa base ai due membri: si uguagliano gli argomenti.
\begin{enumerate}
    \item $x=16$ (C.E. $x>0$ rispettata).
    \item $x+1=4 \Rightarrow x=3$ (C.E. $x>-1$).
    \item $x-2=3 \Rightarrow x=5$ (C.E. $x>2$).
    \item $x-1=3 \Rightarrow x=4$ (C.E. $x>1$).
\end{enumerate}
\end{explanation}
\end{question}

%---------------------------------------------------------------------------------%

\begin{question}[ID=2.4, subtitle={Equazioni logaritmiche con proprietà di somma}, difficulty=media, tags={equazioni logaritmiche, proprietà}]
Risolvi:
\begin{center}\begin{minipage}{0.92\linewidth}
\begin{multicols}{2}
\begin{enumerate}[leftmargin=3.2cm]
\item $\log_2(x) + \log_2(x) = 4$
\item $\log_3(x) + \log_3(2) = \log_3(18)$
\item $\log(x) + \log(x) = \log(100)$
\item $\log_5(x^2) = 2$
\end{enumerate}
\end{multicols}
\end{minipage}\end{center}
\begin{solution}
\begin{enumerate}
    \item $x=4$
    \item $x=9$
    \item $x=10$
    \item $x=\pm5$
\end{enumerate}
\end{solution}
\begin{explanation}
Prima si scrivono le \textbf{condizioni di esistenza} (argomenti positivi); poi si usa la definizione $\log_a A=b \iff A=a^b$, oppure il confronto $\log_aA=\log_aB \iff A=B$; infine si controllano le C.E.
Proprietà dei logaritmi (argomenti positivi): \ $\log_a(xy)=\log_ax+\log_ay$; \ $\log_a\dfrac xy=\log_ax-\log_ay$; \ $\log_a x^n=n\log_ax$; \ $\log_a a^n=n$.
\begin{enumerate}
    \item C.E. $x>0$. \ $2\log_2x=4 \Rightarrow \log_2x=2 \Rightarrow x=4$.
    \item C.E. $x>0$. \ $\log_3(2x)=\log_318 \Rightarrow 2x=18 \Rightarrow x=9$.
    \item C.E. $x>0$. \ $2\log x=2 \Rightarrow \log x=1 \Rightarrow x=10$.
    \item C.E. $x^2>0$, cioè $x\neq0$. \ $x^2=5^2=25 \Rightarrow x=\pm5$, entrambe accettabili. \\
    (Attenzione: scrivere $2\log_5x$ farebbe perdere $x=-5$, perché richiederebbe $x>0$.)
\end{enumerate}
\end{explanation}
\end{question}

%---------------------------------------------------------------------------------%

\begin{question}[ID=2.5, subtitle={Equazioni esponenziali risolte con i logaritmi}, difficulty=media, tags={equazioni esponenziali, logaritmi}]
Risolvi le seguenti equazioni esponenziali utilizzando i logaritmi:
\begin{center}\begin{minipage}{0.92\linewidth}
\begin{multicols}{2}
\begin{enumerate}[leftmargin=3.2cm]
\item $2^x = 8$
\item $3^x = 20$
\item $10^x = 7$
\item $5^{x-1} = 12$
\end{enumerate}
\end{multicols}
\end{minipage}\end{center}
\begin{solution}
\begin{enumerate}
    \item $x=3$
    \item $x=\log_320\approx2{,}73$
    \item $x=\log7\approx0{,}85$
    \item $x=1+\log_512\approx2{,}54$
\end{enumerate}
\end{solution}
\begin{explanation}
Se i due membri non si possono scrivere con la stessa base, si applica il logaritmo: $a^x=b \iff x=\log_ab$. Per calcolarlo si usa il cambio di base $\log_ab=\dfrac{\log b}{\log a}$.
\begin{enumerate}
    \item $x=\log_28=3$.
    \item $x=\log_320=\dfrac{\log20}{\log3}\approx\dfrac{1{,}301}{0{,}477}\approx2{,}73$.
    \item $x=\log_{10}7\approx0{,}85$.
    \item $x-1=\log_512 \Rightarrow x=1+\log_512\approx1+1{,}54$.
\end{enumerate}
\end{explanation}
\end{question}

%====================================================%
% SEZIONE 3: Disequazioni
%====================================================%
\begin{question}[ID=3.1, subtitle={Disequazioni logaritmiche elementari}, difficulty=facile, tags={disequazioni logaritmiche, definizione}]
Risolvi le seguenti disequazioni logaritmiche:
\begin{center}\begin{minipage}{0.92\linewidth}
\begin{multicols}{2}
\begin{enumerate}[leftmargin=3.2cm]
\item $\log_3 x > 2$
\item $\log_{\frac{1}{3}} x > 1 \color{gray}$
\item $\color{gray}\log(2x+1) > \log(x-1) \color{gray}$
\item $\color{gray}\ln\!\left( \tfrac{1}{2} x -1 \right) < 0$
\end{enumerate}
\end{multicols}
\end{minipage}\end{center}
\begin{solution}
\begin{enumerate}
    \item $x>9$
    \item $0<x<\tfrac13$
    \item $x>1$
    \item $2<x<4$
\end{enumerate}
\end{solution}
\begin{explanation}
Nelle disequazioni si mettono a sistema le C.E. con la disequazione tra gli argomenti: il verso \textbf{si conserva} se la base è maggiore di $1$, \textbf{si inverte} se è tra $0$ e $1$.
\begin{enumerate}
    \item C.E. $x>0$; base $3>1$: $x>3^2=9$.
    \item C.E. $x>0$; base $\tfrac13<1$: il verso si inverte, $x<\left(\tfrac13\right)^1$. \ Con le C.E.: $0<x<\tfrac13$.
    \item \[ \begin{cases} 2x+1>0\\ x-1>0\\ 2x+1>x-1 \end{cases} \;\Rightarrow\; \begin{cases} x>-\dfrac12\\ x>1\\ x>-2 \end{cases} \;\Rightarrow\; x>1 \]
    \item $\ln t<0 \iff 0<t<1$:
    \[ 0<\dfrac12x-1<1 \;\Rightarrow\; 1<\dfrac12x<2 \;\Rightarrow\; 2<x<4 \]
\end{enumerate}
\end{explanation}
\end{question}

%---------------------------------------------------------------------------------%

\begin{question}[ID=3.2, subtitle={Disequazioni logaritmiche elementari (variante)}, difficulty=facile, tags={disequazioni logaritmiche, definizione}]
Risolvi le seguenti disequazioni logaritmiche:
\begin{center}\begin{minipage}{0.92\linewidth}
\begin{multicols}{2}
\begin{enumerate}[leftmargin=3.2cm]
\item $\log_2 x < 3$
\item $\log_{\frac{1}{2}} x \geq 2 \color{gray}$
\item $\color{gray}\log(3x-1) \leq \log(x+2)$
\item $\ln(x-2) > 0$
\end{enumerate}
\end{multicols}
\end{minipage}\end{center}
\begin{solution}
\begin{enumerate}
    \item $0<x<8$
    \item $0<x\leq\tfrac14$
    \item $\tfrac13<x\leq\tfrac32$
    \item $x>3$
\end{enumerate}
\end{solution}
\begin{explanation}
Nelle disequazioni si mettono a sistema le C.E. con la disequazione tra gli argomenti: il verso \textbf{si conserva} se la base è maggiore di $1$, \textbf{si inverte} se è tra $0$ e $1$.
\begin{enumerate}
    \item C.E. $x>0$; $x<2^3=8$.
    \item Base $\tfrac12<1$: $x\leq\left(\tfrac12\right)^2=\tfrac14$; con $x>0$.
    \item \[ \begin{cases} 3x-1>0\\ x+2>0\\ 3x-1\leq x+2 \end{cases} \;\Rightarrow\; \begin{cases} x>\dfrac13\\ x>-2\\ x\leq\dfrac32 \end{cases} \]
    \item $\ln(x-2)>0 \iff x-2>1 \iff x>3$.
\end{enumerate}
\end{explanation}
\end{question}

%---------------------------------------------------------------------------------%

\begin{question}[ID=3.3, subtitle={Disequazioni logaritmiche per confronto}, difficulty=media, tags={disequazioni logaritmiche, confronto}]
Risolvi le seguenti disequazioni:
\begin{center}\begin{minipage}{0.92\linewidth}
\begin{multicols}{2}
\begin{enumerate}[leftmargin=3.2cm]
\item $\log_2(x) > \log_2(16)$
\item $\log_3(x+1) \le \log_3(4)$
\item $\log(x-2) \ge \log(3)$
\item $\log_4(x-1) < \log_4(3)$
\item $\log_5(2x) > \log_5(10)$
\item $5^x \le 12$
\end{enumerate}
\end{multicols}
\end{minipage}\end{center}
\begin{solution}
\begin{enumerate}
    \item $x>16$
    \item $-1<x\leq3$
    \item $x\geq5$
    \item $1<x<4$
    \item $x>5$
    \item $x\leq\log_512$
\end{enumerate}
\end{solution}
\begin{explanation}
Nelle disequazioni si mettono a sistema le C.E. con la disequazione tra gli argomenti: il verso \textbf{si conserva} se la base è maggiore di $1$, \textbf{si inverte} se è tra $0$ e $1$.
Tutte le basi sono maggiori di $1$: il verso si conserva.
\begin{enumerate}
    \item $x>16$ (che implica $x>0$).
    \item $x+1\leq4$ con $x+1>0$: \ $-1<x\leq3$.
    \item $x-2\geq3$: \ $x\geq5$.
    \item $x-1<3$ con $x-1>0$: \ $1<x<4$.
    \item $2x>10$: \ $x>5$.
    \item $5^x\leq12 \Rightarrow x\leq\log_512$ (applico $\log_5$, funzione crescente).
\end{enumerate}
\end{explanation}
\end{question}

%---------------------------------------------------------------------------------%

\begin{question}[ID=3.4, subtitle={Disequazioni logaritmiche con proprietà di somma}, difficulty=media, tags={disequazioni logaritmiche, proprietà}]
Risolvi:
\begin{center}\begin{minipage}{0.92\linewidth}
\begin{multicols}{2}
\begin{enumerate}[leftmargin=3.2cm]
\item $\log_2(x) + \log_2(x) > 4$
\item $\log_3(x) + \log_3(2) \le \log_3(18)$
\item $\log(x) + \log(x) \ge \log(100)$
\item $\log_5(x^2) < 2$
\item $\log_7(x) + \log_7(x) \ge \log_7(49)$
\item $2^x > 5$
\end{enumerate}
\end{multicols}
\end{minipage}\end{center}
\begin{solution}
\begin{enumerate}
    \item $x>4$
    \item $0<x\leq9$
    \item $x\geq10$
    \item $-5<x<5$, \ $x\neq0$
    \item $x\geq7$
    \item $x>\log_25$
\end{enumerate}
\end{solution}
\begin{explanation}
Nelle disequazioni si mettono a sistema le C.E. con la disequazione tra gli argomenti: il verso \textbf{si conserva} se la base è maggiore di $1$, \textbf{si inverte} se è tra $0$ e $1$.
Proprietà dei logaritmi (argomenti positivi): \ $\log_a(xy)=\log_ax+\log_ay$; \ $\log_a\dfrac xy=\log_ax-\log_ay$; \ $\log_a x^n=n\log_ax$; \ $\log_a a^n=n$.
\begin{enumerate}
    \item C.E. $x>0$: \ $2\log_2x>4 \Rightarrow \log_2x>2 \Rightarrow x>4$.
    \item C.E. $x>0$: \ $\log_3(2x)\leq\log_318 \Rightarrow 2x\leq18 \Rightarrow x\leq9$.
    \item C.E. $x>0$: \ $\log x^2\geq\log100 \Rightarrow x^2\geq100$; con $x>0$: $x\geq10$.
    \item C.E. $x\neq0$: \ $x^2<25 \Rightarrow -5<x<5$, escluso $0$.
    \item C.E. $x>0$: \ $x^2\geq49 \Rightarrow x\geq7$.
    \item $x>\log_25$.
\end{enumerate}
\end{explanation}
\end{question}

%---------------------------------------------------------------------------------%

\begin{question}[ID=3.5, subtitle={Disequazioni esponenziali risolte con i logaritmi}, difficulty=media, tags={disequazioni esponenziali, logaritmi}]
Risolvi le seguenti disequazioni esponenziali:
\begin{center}\begin{minipage}{0.92\linewidth}
\begin{multicols}{2}
\begin{enumerate}[leftmargin=3.2cm]
\item $2^x > 8$
\item $3^x \le 20$
\item $10^x \ge 7$
\item $5^{x-1} < 12$
\item $e^x > 15$
\item $4^{2x} \le 10$
\end{enumerate}
\end{multicols}
\end{minipage}\end{center}
\begin{solution}
\begin{enumerate}
    \item $x>3$
    \item $x\leq\log_320$
    \item $x\geq\log7$
    \item $x<1+\log_512$
    \item $x>\ln15$
    \item $x\leq\tfrac12\log_410=\log_{16}10\approx0{,}83$
\end{enumerate}
\end{solution}
\begin{explanation}
Si applica ai due membri il logaritmo nella base della potenza: le basi sono tutte maggiori di $1$, quindi il verso si conserva.
\begin{enumerate}
    \item $2^x>2^3 \Rightarrow x>3$.
    \item $x\leq\log_320$.
    \item $x\geq\log_{10}7$.
    \item $x-1<\log_512 \Rightarrow x<1+\log_512$.
    \item $x>\ln15$ (logaritmo in base $e$).
    \item \[\begin{aligned}
    4^{2x} &\leq 10\\
    2x &\leq \log_410\\
    x &\leq \dfrac12\log_410=\dfrac{\log10}{2\log4}=\dfrac{1}{4\log2}\approx0{,}83 && \textcolor{gray}{\text{cambio di base, $\log4=2\log2$}}
    \end{aligned}\]
    (Equivalentemente $4^{2x}=16^x$, quindi $x\leq\log_{16}10$.)
\end{enumerate}
\end{explanation}
\end{question}
